Kolekcija materijala sa vežbi za kurs \href{http://www.verifikacijasoftvera.matf.bg.ac.rs/}{Verifikacija softvera} na Matematičkom fakultetu, Univerziteta u Beogradu.

Materijali prate praktičan tok rada: od dijagnoziranja simptoma i debagovanja, preko testiranja i profajliranja, do simboličkog izvršavanja, proveravanja modela i formalne verifikacije. Svaka tema je organizovana kroz kratko objašnjenje, komande, očekivane nalaze i male primere koji se nalaze uz odgovarajuće poglavlje.

Pokrivene su sledeće celine:

\begin{itemize}
\tightlist
\item
  \textbf{Dijagnoziranje problema}: osnovni Linux alati, posmatranje procesa i analiza ponašanja programa metodom crne kutije.
\item
  \textbf{Debagovanje}: \texttt{gdb}, QtCreator, core dump analiza, watchpoint-i, niti, prikačivanje na proces, Java JDWP, agenti i native/AOT izvršive verzije.
\item
  \textbf{Refaktorisanje}: prepoznavanje lošeg dizajna, izdvajanje odgovornosti i primeri kroz C++, Java i C\# kod.
\item
  \textbf{Testiranje jedinica koda}: QtTest, JUnit, xUnit, NUnit, pokrivenost koda, \texttt{lcov}, JaCoCo i .NET alati za pokrivenost.
\item
  \textbf{Testiranje pomoću objekata imitatora}: ručni imitatori, Mockito, Moq i testiranje koda koji zavisi od servisa, vremena ili baze podataka.
\item
  \textbf{Fuzz testiranje}: libFuzzer, tmux tokovi rada, Jazzer za Java kod i AFL++ na realnijem C/C++ primeru.
\item
  \textbf{Profajliranje instrumentacijom}: Valgrind alati kao što su Memcheck, Massif, DHAT, Cachegrind, Callgrind, Helgrind i DRD, kao i \texttt{gprof} i LLVM XRay.
\item
  \textbf{Profajliranje uzorkovanjem}: \texttt{perf}, flame graph tok rada, Intel VTune Profiler i Java async-profiler.
\item
  \textbf{eBPF}: model izvršavanja, verifier, mape, \texttt{bpftrace}, BCC alati, sistemsko profajliranje i ograničenja dozvola.
\item
  \textbf{KLEE}: simboličko izvršavanje, generisanje testova, pronalaženje grešaka, rad sa bibliotekama i heuristike pretrage.
\item
  \textbf{CBMC}: proveravanje svojstava C programa, odmotavanje petlji, pretpostavke, postuslovi, harness funkcije i konkurentni primeri.
\item
  \textbf{Clang alati}: Clang statički analizator, \texttt{clang-tidy}, sanitizeri, pojedinačni primeri grešaka i analiza projekata.
\item
  \textbf{Dafny}: specifikacije, pre/postuslovi, invarijante petlji i dokazi korektnosti za male algoritme.
\item
  \textbf{Instalacije}: kratke beleške za alate i pakete koji su potrebni za lokalno pokretanje vežbi.
\end{itemize}
